\documentclass[11pt]{article}
\usepackage[margin=1.1in]{geometry}
\usepackage{amsmath, amssymb, booktabs}
\usepackage[colorlinks=true, linkcolor=blue]{hyperref}

\newcommand{\E}{\mathbb{E}}
\newcommand{\Post}{\mathrm{Post}}

\title{Institution-Level Treatment Effects:\\ Specification, Identification, and the Levels-vs-Slopes Distinction}
\author{}
\date{\today}

\begin{document}
\maketitle

\section{Setup and notation}

PI-by-year panel, 2010--2019. PI $i$ belongs to exactly one institution
$j(i)$ for the entire panel (no movers, by sample construction). The 2014
procurement shock hits PIs with intensity $E_i$ (imputed exposure, fixed over
time); $S_i$ is the PI's market-spend share (the partialling-out control from
the main reduced form). Define
\[
\Post_t = \mathbf{1}[t \ge 2014], \qquad
Z_{it} = E_i \cdot \Post_t, \qquad
Z^{S}_{it} = S_i \cdot \Post_t .
\]
Outcome $Y_{it}$ is the last-author publication count; all models are PPML
(Poisson pseudo-likelihood), so coefficients read as log points.

\section{Three distinct institution-level objects}

It is worth being pedantic about three objects that are easy to conflate.
Write the conditional mean in an encompassing form:
\begin{equation}
\E[Y_{it}] \;=\; \exp\Big(
\underbrace{\alpha_i}_{\text{PI level}}
+ \underbrace{\delta_t}_{\text{year}}
+ \underbrace{\phi_{j(i)}}_{\substack{\text{institution}\\ \text{LEVEL}}}
+ \underbrace{\pi_{j(i)}\,\Post_t}_{\substack{\text{institution}\\ \text{response INTERCEPT}}}
+ \underbrace{\beta_{j(i)}\, Z_{it}}_{\substack{\text{institution}\\ \text{response SLOPE}}}
+ \gamma_{j(i)}\, Z^{S}_{it}
\Big).
\label{eq:encompassing}
\end{equation}

To fix language, consider institution $j$'s \emph{response line}: the
post-minus-pre change in log expected output for one of its PIs with exposure
$E$,
\begin{equation}
\Delta_j(E) \;=\; \pi_j + \beta_j \, E .
\label{eq:responseline}
\end{equation}

\begin{center}
\begin{tabular}{@{}llll@{}}
\toprule
Object & Name used here & What it measures & Identified w/o movers? \\
\midrule
$\phi_j$ & institution \textbf{level} & permanent output difference & \textbf{No} \\
$\pi_j$ & response \textbf{intercept} ($te_j$) & post-2014 shift common to all of $j$'s PIs & \textbf{Yes} \\
$\beta_j$ & response \textbf{slope} (gradient) & how the shift scales with exposure within $j$ & Yes (needs $\ge 2$ PIs w/ distinct $E$) \\
\bottomrule
\end{tabular}
\end{center}

\paragraph{``Levels'' vs.\ ``slopes,'' precisely.}
Throughout this note, \emph{level} refers only to $\phi_j$: a time-invariant
institution effect on output. \emph{Intercept} ($\pi_j$, our $te_j$) is a
level shift \emph{of the response} --- constant across the institution's PIs
but occurring only after 2014. \emph{Slope} ($\beta_j$) is the gradient of
the response in exposure, within the institution. The earlier draft of this
analysis estimated institution-specific $\beta_j$; the current version
estimates institution-specific $\pi_j$ with a pooled slope. They answer
different questions (Section~\ref{sec:gradient}).

\section{Why institution levels $\phi_j$ are not identified}

In an AKM design, person and firm effects are separated by movers: when a
worker switches firms, the change in her outcome identifies the firm
contrast. Here nobody moves, so the PI fixed effects are a strict refinement
of institution fixed effects: the institution dummy equals the sum of its
PIs' dummies. Adding a constant $c$ to $\phi_j$ and subtracting $c$ from
$\alpha_i$ for every $i$ with $j(i)=j$ leaves the likelihood unchanged. The
data cannot distinguish a productive institution from an institution staffed
by productive people. Consequently $\phi_j$ is dropped from all estimated
models (absorbed into $\alpha_i$), and \emph{nothing in this analysis speaks
to permanent institutional quality differences}.

\section{The estimated specification: institution response intercepts}
\label{sec:main}

The implemented model (\texttt{inst\_fe\_diag.do}, Stage 1) restricts the
slope to be common across institutions ($\beta_j \equiv b$, $\gamma_j \equiv
g$) and saturates the intercepts:
\begin{equation}
\E[Y_{it}] \;=\; \exp\Big( \alpha_i + \delta_t + b\, Z_{it} + g\, Z^{S}_{it}
+ \sum_{j \ne j_0} te_j \, \mathbf{1}[j(i)=j] \cdot \Post_t \Big),
\label{eq:main}
\end{equation}
estimated by \texttt{ppmlhdfe} with standard errors clustered by PI.

\paragraph{Interpretation.}
Take two PIs with the \emph{same} exposure $E$, one at institution $j$, one
at institution $k$. Under \eqref{eq:main} their expected post-minus-pre
log-output changes differ by exactly $te_j - te_k$: the common exposure
response $b E$ and the common year path $\delta_t$ cancel. So $te_j$ answers:
\emph{conditional on identical exposure, how does being at institution $j$
change the output response to the shock period?} This is the feasible analog
of the AKM institution effect --- but for the \emph{change} rather than the
level.

\paragraph{Why $te_j$ is identified without movers.}
The trick is that $te_j$ multiplies $\Post_t$, which varies \emph{within} PI
over time. A PI fixed effect $\alpha_i$ is constant over the panel; it
absorbs ``PI $i$ is productive,'' but cannot absorb ``PI $i$'s output fell
after 2014.'' Each PI's pre-2014 years therefore act as her own control, and
$te_j$ is (up to the common terms being netted out) an average of
\emph{within-PI pre/post changes} across the PIs of institution $j$. No
comparison requires observing the same person at two institutions. The
person-vs-place confound that destroys $\phi_j$ never arises, because the
person's level is differenced out by her own history. Even a single-PI
institution has an identified (if noisy) $te_j$.

\paragraph{Normalization.}
Summing the interactions over all institutions gives
$\sum_j \mathbf{1}[j(i)=j]\Post_t = \Post_t$, which lies in the span of the
year effects. Hence one institution $j_0$ must be omitted and only
\emph{relative} intercepts are identified --- the same normalization as AKM
firm effects; it is a bookkeeping issue, not an identification failure. We
omit the largest institution (Johns Hopkins, 464 PIs) so the reference is
maximally precise, then re-center all $\hat{te}_j$ at their precision-weighted
mean ($1/\hat{se}_j^2$ weights). Reported effects are therefore
\emph{deviations from the average institution}. Re-centering also nets out
the reference institution's estimation error, which is common to every
$\hat{te}_j$. All second-stage slopes and variance decompositions are
invariant to these normalization choices; only the constant moves.

\paragraph{What $te_j$ does and does not contain.}
$te_j$ absorbs \emph{any} institution-wide change in output coinciding with
the post period --- shock transmission, but also e.g.\ a 2015 building
opening or leadership change. Two facts tie the estimated heterogeneity back
to the shock: (i) institutions with higher mean exposure have systematically
more negative $\hat{te}_j$ (the ``dose'' gradient in the second stage), and
(ii) the buffer covariates (endowment etc.)\ explain the between-institution
variance in the predicted direction. Additionally, without movers, a causal
institution effect cannot be separated from \emph{sorting}: institutions may
differ in the (unobserved) resilience of the PIs they attract. $te_j$ is
``the effect of being at institution $j$'' in the descriptive sense.

\section{The earlier specification: institution response slopes}
\label{sec:gradient}

An earlier version of the analysis instead saturated the \emph{slopes}:
\begin{equation}
\E[Y_{it}] \;=\; \exp\Big( \alpha_i + \lambda_{j(i),t}
+ \beta_j \, Z_{it} + \gamma_j \, Z^{S}_{it} \Big),
\label{eq:gradient}
\end{equation}
where $\lambda_{jt}$ are institution$\times$year fixed effects. Two points
about this design:

\begin{enumerate}
\item \textbf{What $\beta_j$ measures.} $\beta_j$ is the slope of the
response line \eqref{eq:responseline} within institution $j$: among PIs
\emph{at the same institution}, how much worse is the post-2014 outcome per
unit of exposure. Cross-institution heterogeneity in $\beta_j$ asks: is the
exposure \emph{gradient} steeper at some institutions than others?

\item \textbf{Why $\lambda_{jt}$ (or at least $\pi_j \Post_t$) must be
included.} If \eqref{eq:gradient} used only common year effects $\delta_t$,
an institution-wide post shock ($\pi_j$-type variation) would have no
dedicated regressor and would be projected onto the only institution-specific
regressors available --- the exposure interactions --- contaminating
$\hat\beta_j$. With $\lambda_{jt}$ absorbed, every institution-level shift in
every year is soaked up nonparametrically, and $\beta_j$ is a pure
within-institution, cross-PI comparison. The cost is that $\lambda_{jt}$
also absorbs $\pi_j$ itself: \emph{the slope specification is mechanically
silent about intercept heterogeneity, and vice versa.}
\end{enumerate}

\paragraph{Identification requirements differ.}
$\beta_j$ is identified only if institution $j$ has at least two PIs with
\emph{different} exposure values (with exactly two, $\beta_j$ and $\gamma_j$
are collinear and only one is estimable). Single-PI institutions have no
within-institution exposure contrast at all; in the joint estimation their
unidentified columns degrade the full variance matrix --- the practical
failure that motivated the redesign. The intercept specification
\eqref{eq:main} has no such requirement.

\paragraph{The two objects are complements.}
Neither specification can see the other's effect:
\begin{center}
\begin{tabular}{@{}lcc@{}}
\toprule
True institutional behavior & Appears in $te_j$? & Appears in $\beta_j$? \\
\midrule
Uniform cushion (helps all PIs equally) & Yes & No \\
Targeted cushion (protects exposed PIs) & Partially$^{\dagger}$ & Yes (flatter $\beta_j$) \\
Unrelated inst-wide 2014 change & Yes (confound) & No \\
\bottomrule
\end{tabular}

{\footnotesize $^{\dagger}$a targeted cushion raises the institution's
\emph{average} response and so moves $te_j$ toward zero, by an amount
depending on the institution's exposure distribution.}
\end{center}
A buffering story in which well-endowed institutions cushion the shock
predicts \emph{both} a less-negative intercept ($te_j$) and a flatter slope
($\beta_j$) at buffered institutions. The current pipeline tests the
intercept margin; the slope margin is a natural companion exercise using
\eqref{eq:gradient}.

\section{Estimation and inference details}

\begin{itemize}
\item \textbf{Stage 1.} Equation \eqref{eq:main} by \texttt{ppmlhdfe};
institution$\times$post dummies are constructed manually (Stata's factor
machinery silently re-bases interactions, which in one run moved the
reference institution without notice). Clustering is at the PI level:
institution-level clustering is \emph{mechanically degenerate} for
institution-specific coefficients, because each $te_j$'s score has support in
exactly one cluster, driving its cluster-robust variance to zero. PI
clustering understates any residual cross-PI dependence within institutions;
first-stage SEs are therefore screening-grade, with definitive inference
deferred to an institution-block bootstrap of the full pipeline.
\item \textbf{Heterogeneity test.} Cochran's $Q=\sum_j
\hat{se}_j^{-2}(\hat{te}_j - \overline{te})^2 \sim \chi^2_{J-1}$ under
homogeneity (current run: $Q=161.5$, $df=127$, $p=0.021$).
\item \textbf{Stage 2.} Fay--Herriot random-effects meta-regressions of
$\hat{te}_j$ on institution characteristics (REML, Knapp--Hartung SEs),
weighting each institution by $1/(\hat\tau^2 + \hat{se}_j^2)$; a univariate
screen with BH/Bonferroni corrections, then a joint ``horse race'' of
endowment, nonprofit funding, and funding-source diversification, controlling
for institution mean exposure (dose).
\item \textbf{Empirical Bayes.} $\hat{te}^{EB}_j = \bar{te} +
\hat\lambda_j(\hat{te}_j - \bar{te})$ with $\hat\lambda_j = \hat\tau^2 /
(\hat\tau^2 + \hat{se}_j^2)$. Because most of the raw cross-institution
spread is sampling noise ($I^2 \approx 20$--$25\%$), shrinkage is strong;
the EB spread is a lower bound on true heterogeneity given that first-stage
SEs are, if anything, understated.
\item \textbf{Cross-check.} A one-step pooled PPML replacing the estimated
$\hat{te}_j$ with direct $\Post_t \times W_j$ interactions reproduces the
horse-race coefficients in sign and magnitude for endowment and
diversification, validating that the two-step projection is not an artifact
of estimation-then-regression.
\end{itemize}

\end{document}
